% =====================================================================
% _pgfpreamble.tex — shared pgfplots/tikz setup for the TeX-auditable
% figure set. Generated by edit-harness/experiments/make_figures_tex.py.
% \input this once in the main.tex preamble (after \usepackage{pgfplots}).
%
% Palette: the dataviz-skill validated reference categorical instance, used
% in fixed slot order (blue, aqua, yellow, green, violet, red, magenta,
% orange). Sequential = blue ramp; diverging = blue(+)/red(-); status crit.
% Colorblind-safe (worst adjacent light-mode CVD dE 24.2, well past >=12).
% =====================================================================
\definecolor{vizblue}{HTML}{2A78D6}    % slot 1
\definecolor{vizaqua}{HTML}{1BAF7A}    % slot 2
\definecolor{vizyellow}{HTML}{EDA100}  % slot 3
\definecolor{vizgreen}{HTML}{008300}   % slot 4
\definecolor{vizviolet}{HTML}{4A3AA7}  % slot 5
\definecolor{vizred}{HTML}{E34948}     % slot 6
\definecolor{vizmagenta}{HTML}{E87BA4} % slot 7
\definecolor{vizorange}{HTML}{EB6834}  % slot 8
\definecolor{vizseqblue}{HTML}{2A78D6} % sequential ramp step 450
\definecolor{vizink}{HTML}{0B0B0B}     % primary ink (text token)
\definecolor{vizmuted}{HTML}{898781}   % axis / muted labels
\definecolor{vizgrid}{HTML}{E1E0D9}    % hairline gridline
\definecolor{vizcrit}{HTML}{D03B3B}    % status: critical (reference lines)

% pgfplots draws bar-plot (ybar/xbar) legend keys as a two-bar pictogram by
% default -- correct for a single multi-bar series, but reads as "two colors
% per entry" when each addlegendentry already names one color (e.g. Fig. 4b
% by-construction/holdout, Fig. 6b raw/DC). Override to the same one-rectangle
% swatch used by line-plot legends, applied globally so every ybar figure
% in this set is consistent.
\pgfplotsset{
  /pgfplots/bar legend/.style={
    area legend,
    legend image code/.code={
      \draw[##1] (0cm,-0.1cm) rectangle (0.6cm,0.1cm);
    },
  },
}
% Base axis style: thin marks, recessive grid/axes, small ink-token labels.
\pgfplotsset{
  compat=1.18,
  vizbase/.style={
    tick align=outside, tick pos=left,
    axis line style={vizmuted, line width=0.4pt},
    x tick label style={font=\scriptsize},
    y tick label style={font=\scriptsize},
    label style={font=\scriptsize, text=vizink},
    title style={font=\footnotesize\bfseries, text=vizink, align=left,
                 at={(0.0,1.03)}, anchor=south west},
    legend style={font=\scriptsize, draw=none, fill=none, text=vizink,
                  inner sep=1pt, row sep=-1pt},
    legend cell align=left,
    ymajorgrids=true,
    major grid style={vizgrid, line width=0.3pt},
    every axis plot/.append style={line width=0.9pt},
    error bars/error bar style={vizmuted, line width=0.5pt},
    error bars/error mark options={rotate=90, mark size=1.3pt, vizmuted},
    clip=false,
  },
}
% Central small-multiples style: every paper panel uses this so the whole figure
% set lays out at REAL dimensions (no \resizebox). Sizes are set explicitly per
% figure (width/height budgeted from \columnwidth=219pt / \textwidth=455pt) so
% pgfplots reports genuine overfulls. paperpanel = vizbase + tighter y-labels +
% ONE canonical legend placement: a horizontal row ABOVE the panel, clear of the
% title (which vizbase pins at y=1.03, left) and never on the data. \tiny keeps
% wide 2-entry legends from overhanging the narrow multiples.
\pgfplotsset{
  paperpanel/.style={
    vizbase,
    ylabel near ticks,
    legend style={at={(0.5,1.30)}, anchor=south, legend columns=-1,
                  draw=none, fill=none, font=\tiny, inner sep=1pt,
                  /tikz/every even column/.append style={column sep=6pt}},
    legend cell align=left,
  },
}
% direct-label node styles (tikz, not pgfplots keys): values/labels wear ink
% tokens, never the series color.
\tikzset{
  vizlabel/.style={font=\scriptsize, text=vizink, inner sep=1pt},
  vizlabelm/.style={font=\scriptsize, text=vizmuted, inner sep=1pt},
}
